\documentclass[11pt]{article}
\usepackage[a4paper,margin=25mm]{geometry}
\usepackage{amsmath,amssymb,amsthm}
\usepackage{longtable}
\usepackage[hidelinks]{hyperref}
\hypersetup{pdftitle={E1774175\_6: x\textasciicircum 4+xy-x+y\textasciicircum 3+2y-1=0 has no integer solutions}}
\newtheorem{theorem}{Theorem}
\newtheorem{lemma}[theorem]{Lemma}
% keep the space around binary operators in inline formulas
\medmuskip=4mu plus 2mu\relax
\emergencystretch=1.5em
\tolerance=400
\allowdisplaybreaks
\newcommand{\Z}{\mathbb{Z}}
\newcommand{\Q}{\mathbb{Q}}
% Legendre symbols: one fixed size in running text, one in displays
\newcommand{\leg}[2]{\mathchoice{\biggl(\dfrac{#1}{#2}\biggr)}{(#1/#2)}{(#1/#2)}{(#1/#2)}}
\title{\textbf{E1774175\_6}\\[4pt] {\Large The equation $x^{4}+xy-x+y^{3}+2y-1=0$ has no integer solutions}}
\author{}
\date{}
\begin{document}
\maketitle
\vspace{-8ex}
\begin{center}
Size $h=35$, length $l=12$
\end{center}

\begin{theorem}\label{thm:main}
The equation
\begin{equation}\label{eq:main}
x^{4}+xy-x+y^{3}+2y-1=0
\end{equation}
has no integer solutions.
\end{theorem}

The proof uses two polynomial identities of the form $H_kT_k=U_k^2+d_kV_k^2$ on the curve \eqref{eq:main}, where $d_{1}=-2H_{2}$ and $d_{2}=-2H_{1}$, and the product formula for the Hilbert symbol of one pair of values. Apart from polynomial identities, which can be checked by expanding both sides, and finite checks of congruences, the proof uses only basic properties of the Legendre symbol and of the Hilbert symbol.

\section{Legendre symbols and Hilbert symbols}

For a prime $p$ and a non-zero integer $N$, we write $v_p(N)$ for the exponent of $p$ in the prime factorization of $N$. For an integer $a$ and an odd prime $p$, the Legendre symbol $\leg{a}{p}$ is $0$ if $p\mid a$, $1$ if $a$ is a non-zero square modulo $p$, and $-1$ otherwise. It is multiplicative in $a$, and it depends only on $a$ modulo $p$ \cite[Chapter~5]{IR}.

For a prime $p$, let $\Q_p$ be the field of $p$-adic numbers, and put $\Q_\infty=\mathbb R$. For $v$ a prime or $v=\infty$, and non-zero $a,b\in\Q$, the Hilbert symbol $(a,b)_v$ is $1$ if the equation $z^2=ar^2+bs^2$ has a solution $(r,s,z)\neq(0,0,0)$ in $\Q_v$, and $-1$ otherwise. In particular, $(a,b)_v=1$ for every $v$ if this equation has a non-zero solution in $\Q$. For an odd integer $t$, we put $\varepsilon(t)=(t-1)/2$ and $\omega(t)=(t^2-1)/8$. We use the following facts \cite[Chapter~III, Proposition~2 and Theorems~1 and~3]{Se}.

\begin{theorem}\label{thm:hilbert}
Let $a$, $a'$ and $b$ be non-zero integers, and let $v$ be a prime or $v=\infty$.
\begin{itemize}
\item[(i)] $(a,b)_\infty=-1$ if and only if $a<0$ and $b<0$.
\item[(ii)] Let $p$ be an odd prime, and write $a=p^\alpha u$ and $b=p^\beta w$ with $p\nmid uw$. Then $(a,b)_p=(-1)^{\alpha\beta\varepsilon(p)}\leg{u}{p}^\beta\leg{w}{p}^\alpha$.
\item[(iii)] Write $a=2^\alpha u$ and $b=2^\beta w$ with $u$ and $w$ odd. Then $(a,b)_2=(-1)^{\varepsilon(u)\varepsilon(w)+\alpha\omega(w)+\beta\omega(u)}$.
\item[(iv)] $(a,b)_v=1$ for all but finitely many $v$, and $\prod_v(a,b)_v=1$, the product over all primes $v$ and $v=\infty$.
\item[(v)] $(a,b)_v=(b,a)_v$ and $(aa',b)_v=(a,b)_v(a',b)_v$.
\end{itemize}
\end{theorem}

\begin{lemma}\label{lem:good}
Let $p$ be an odd prime, and let $a,b,r,c,t$ be integers with $a\neq0$, $b\neq0$ and $at=r^2-bc^2$. If $p\nmid b$, and $p$ does not divide both $r$ and $c$, then $(a,b)_p=1$.
\end{lemma}

\begin{proof}
If $p\nmid a$, this is Theorem~\ref{thm:hilbert}(ii) with $\alpha=\beta=0$. If $p\mid a$, then $r^2\equiv bc^2\pmod p$. If $p\mid c$, then $p\mid r$, which is excluded. Hence $p\nmid c$, so $b\equiv(rc')^2\pmod p$, where $cc'\equiv1\pmod p$, and $\leg{b}{p}=1$ as $p\nmid b$. Theorem~\ref{thm:hilbert}(ii) with $\beta=0$ gives $(a,b)_p=\leg{b}{p}^{v_p(a)}=1$.
\end{proof}

\begin{lemma}\label{lem:cofactor}
Let $a,b,r,c,t$ be integers with $at=r^2-bc^2\neq0$ and $b\neq0$. Then $(a,b)_v=(t,b)_v$ for every $v$.
\end{lemma}

\begin{proof}
The triple $(r,s,z)=(1,c,r)$ is a non-zero rational solution of $z^2=(at)r^2+bs^2$, so $(at,b)_v=1$. By Theorem~\ref{thm:hilbert}(v), $(a,b)_v(t,b)_v=(at,b)_v=1$, and both symbols are $\pm1$.
\end{proof}

\begin{lemma}\label{lem:square}
Let $a$ and $b$ be non-zero integers. If $p$ is an odd prime with $p\nmid a$ and $\leg{a}{p}=1$, then $(a,b)_p=1$. If $a\equiv1\pmod8$, then $(a,b)_2=1$.
\end{lemma}

\begin{proof}
The first claim is Theorem~\ref{thm:hilbert}(ii) with $\alpha=0$. For the second, $\varepsilon(a)$ and $\omega(a)$ are even, and Theorem~\ref{thm:hilbert}(iii) with $\alpha=0$ gives $(a,b)_2=1$.
\end{proof}


\begin{lemma}\label{lem:amgm}
Let $a,b,c$ be real numbers with $a>0$, $c>0$ and $256ac^3>27b^4$. Then $at^4+bt+c>0$ for every real number $t$.
\end{lemma}

\begin{proof}
For $t=0$ this is clear. For $t\neq0$, the inequality between the arithmetic and geometric means of $at^4$, $c/3$, $c/3$ and $c/3$ gives $at^4+c\geq4\,(at^4c^3/27)^{1/4}=(256ac^3/27)^{1/4}|t|>|b|\,|t|\geq-bt$.
\end{proof}

\section{The auxiliary polynomials}

Put $F(x,y)=x^{4}+xy-x+y^{3}+2y-1$, so that \eqref{eq:main} is the equation $F(x,y)=0$. For the first identity, we define
\begin{align*}
H_{1}&=2x^{2}+2xy-2x+y^{2}+4y+6,\\
T_{1}&=36x^{6}-12x^{5}-72x^{4}y-1259x^{4}+72x^{3}y^{2}-432x^{3}y+842x^{3}\\
&\quad+24x^{2}y^{2}+1078x^{2}y+9917x^{2}-1734xy^{2}+17050xy+1490x-3500y^{2}\\
&\quad+6588y+28049,\\
U_{1}&=-6x^{3}y+4x^{3}+x^{2}y+136x^{2}+215xy-64x+201y+402,\\
V_{1}&=19,\\
G_{1}&=72x^{4}+72x^{3}y-96x^{3}+24x^{2}y-2294x^{2}-1734xy+3114x-3500y-7412,\\
\intertext{for the second identity, we define}
H_{2}&=-2y+1,\\
T_{2}&=32x^{3}+32x^{2}y^{2}+16x^{2}y+56x^{2}-16xy^{2}-8xy-64x-30y^{2}-14y-63,\\
U_{2}&=-8x^{3}+2x^{2}+4x+3,\\
V_{2}&=1,\\
G_{2}&=-64x^{2}+32x+60,\\
\intertext{for the auxiliary polynomials $H_{1}$ and $H_{2}$ together, we define}
\lambda_{1}&=-9216x^{3}-32992x^{2}+21520x+151460,\\
\mu_{1}&=-9216x^{4}-4608x^{3}y-53728x^{3}-16496x^{2}y-43496x^{2}+9216xy^{2}\\
&\quad+15368xy+249000x+28384y^{2}+89922y+404649,\\
\nu_{1}&=18432x+56768,\\
\intertext{and, for the signs, we also use}
L_{2}&=y-1,\\
Q_{2}&=y^{3}+2y-1,\\
\Delta_{2}&=256y^{9}+1536y^{7}-768y^{6}+3072y^{5}-3099y^{4}+2924y^{3}-3234y^{2}\\
&\quad+1644y-283.
\end{align*}
Expanding both sides, one checks the identities
\begin{align}
H_{1}T_{1}-U_{1}^2+2H_{2}V_{1}^2&=G_{1}F,\label{eq:norm1}\\
H_{2}T_{2}-U_{2}^2+2H_{1}V_{2}^2&=G_{2}F,\label{eq:norm2}\\
\lambda_{1}H_{1}+\mu_{1}H_{2}+\nu_{1}F&=19^{2}\cdot 59^{2},\label{eq:cross1}\\
H_{1}&=(x+y+2)^2+(x-3)^2-7.\label{eq:sign1}\\
F&=x^4+L_{2}x+Q_{2},\label{eq:sign2}\\
\Delta_{2}&=256Q_{2}^3-27L_{2}^4.\label{eq:sign2b}
\end{align}

\section{Proof of Theorem~\ref{thm:main}}

Suppose that $(x,y)\in\Z^2$ is a solution of \eqref{eq:main}, so that $F(x,y)=0$; from now on the polynomials denote their values at $(x,y)$. Put
\begin{gather*}
A=-d_{2}=2H_{1},\quad B=-d_{1}=2H_{2},\quad T_A=2T_{1},\quad T_B=2T_{2}.
\end{gather*}
By \eqref{eq:norm1}--\eqref{eq:norm2}, after multiplying by squares of the constants $c$ in $d_k=cH_{k'}$ if necessary, we obtain
\begin{align}
AT_A=(2U_{1})^2-B(2V_{1})^2,\qquad BT_B=(2U_{2})^2-A(2V_{2})^2.\label{eq:AB}
\end{align}

\paragraph{Step 1: signs.} \emph{The sign of $H_{1}$.} If $H_{1}\leq0$, then \eqref{eq:sign1} gives $(x+y+2)^2\leq 7$ and $(x-3)^2\leq 7$, so $|x+y+2|\leq2$ and $|x-3|\leq2$. Solving these two linear forms for $x$ and $y$, we find that the only integer points satisfying both inequalities are $(1,-5)$, $(1,-4)$, $(1,-3)$, $(1,-2)$, $(1,-1)$, $(2,-6)$, $(2,-5)$, $(2,-4)$, $(2,-3)$, $(2,-2)$, $(3,-7)$, $(3,-6)$, $(3,-5)$, $(3,-4)$, $(3,-3)$, $(4,-8)$, $(4,-7)$, $(4,-6)$, $(4,-5)$, $(4,-4)$, $(5,-9)$, $(5,-8)$, $(5,-7)$, $(5,-6)$, $(5,-5)$, and none of them is a solution of \eqref{eq:main} with $H_{1}\leq0$. Hence $H_{1}>0$.

\emph{The sign of $H_{2}$.} Suppose that $H_{2}\leq0$. Since $H_{2}=-2y+1$, this means $y\geq1$. By \eqref{eq:sign2}, $F$ is $x^4+L_{2}x+Q_{2}$, where $L_{2}$ and $Q_{2}$ depend only on $y$. If $y\geq1$, write $y=1+z$ with $z\geq0$; expanding $Q_{2}$ and $\Delta_{2}$ as polynomials in $z$ gives non-negative coefficients and positive constant terms, so $Q_{2}>0$ and $\Delta_{2}>0$. Hence $Q_{2}>0$ and $\Delta_{2}>0$, so Lemma~\ref{lem:amgm} with $a=1$, $b=L_{2}$, $c=Q_{2}$ and $t=x$ gives $F(x,y)>0$, a contradiction. Hence $H_{2}>0$.

By these signs, $A>0$ and $B>0$.

Hence Theorem~\ref{thm:hilbert}(i) gives $(A,B)_\infty=1$ (as $A>0$).

\paragraph{Step 2: primes outside $S$.} Let $S=\{2,19,59\}$ be the set of primes dividing $2$, the constant $2$ in the $d_k=cH_{k'}$, $V_{1}=19$ and $V_{2}=1$, and the right-hand side of \eqref{eq:cross1}. Let $p\notin S$ be an odd prime; we show that $(A,B)_p=1$. If $p\nmid AB$, this follows from Theorem~\ref{thm:hilbert}(ii) with $\alpha=\beta=0$. If $p\mid A$, then $p\mid H_{1}$, so $p\nmid H_{2}$ by \eqref{eq:cross1} and $p\nmid B$; as $p\nmid 2V_{1}$, Lemma~\ref{lem:good} with $(a,b,r,c,t)=(A,B,2U_{1},2V_{1},T_A)$ and \eqref{eq:AB} gives $(A,B)_p=1$; if $p\mid B$, then $p\mid H_{2}$, so $p\nmid H_{1}$ by \eqref{eq:cross1} and $p\nmid A$; as $p\nmid 2V_{2}$, Lemma~\ref{lem:good} with $(a,b,r,c,t)=(B,A,2U_{2},2V_{2},T_B)$ and \eqref{eq:AB} gives $(B,A)_p=1$, which equals $(A,B)_p$ by Theorem~\ref{thm:hilbert}(v).

\paragraph{Step 3: the primes in $S$.} By Step~1, $A$ and $B$ are non-zero. If $T_A\neq0$, then \eqref{eq:AB} and Lemma~\ref{lem:cofactor} give $(A,B)_v=(T_A,B)_v$, and if $T_B\neq0$, they give $(A,B)_v=(B,A)_v=(T_B,A)_v=(A,T_B)_v$, using Theorem~\ref{thm:hilbert}(v). For a non-zero integer $X$ and a prime $p$, we write $X\in p^{2i}(\pm)$ or $X\in p^{2i+1}(\pm)$, when $p$ is odd, if $X=p^ju$ with $p\nmid u$, $\leg{u}{p}=\pm1$, and $j$ even or odd, respectively; and we write $X\in2^{2i}\cdot u$ or $X\in2^{2i+1}\cdot u$ if $X=2^ju'$ with $u'\equiv u\pmod 8$ and $j$ even or odd. If $X\not\equiv0\pmod{p^e}$ and $v_p(X)\leq e-1$, or $v_p(X)\leq e-3$ if $p=2$, then this class is determined by $X$ modulo $p^e$. By Theorem~\ref{thm:hilbert}(ii) and~(iii), the classes of $X$ and $Y$ determine $(X,Y)_p$.

\emph{The prime $2$.} Checking the $16^2$ residue pairs modulo $16$, we find that \eqref{eq:main} has exactly 8 solutions modulo $16$, and that at each of them one of the cases listed for $p=2$ in Table~\ref{tab:local} holds. Each case gives the stated value of $(A,B)_{2}$. Hence $(A,B)_{2}=-1$.

\emph{The prime $19$.} Checking the $19^2$ residue pairs modulo $19$, we find that \eqref{eq:main} has exactly 17 solutions modulo $19$, and that at each of them one of the cases listed for $p=19$ in Table~\ref{tab:local} holds. Each case gives the stated value of $(A,B)_{19}$. Hence $(A,B)_{19}=1$.

\emph{The prime $59$.} Checking the $59^2$ residue pairs modulo $59$, we find that \eqref{eq:main} has exactly 59 solutions modulo $59$, and that at each of them one of the cases listed for $p=59$ in Table~\ref{tab:local} holds. Each case gives the stated value of $(A,B)_{59}$. Hence $(A,B)_{59}=1$.

\paragraph{Step 4: contradiction.} By Theorem~\ref{thm:hilbert}(iv), $(A,B)_\infty\prod_p(A,B)_p=1$, the product over all primes, and by Step~2, the factors with $p\notin S$ are $1$. Hence $\prod_{p\in S}(A,B)_p=(A,B)_\infty=1$. But by Step~3, $\prod_{p\in S}(A,B)_p=-1$. This contradiction proves the theorem.
\qed

\begin{thebibliography}{9}
\bibitem{IR} K. Ireland and M. Rosen, \emph{A Classical Introduction to Modern Number Theory}, 2nd ed., Graduate Texts in Mathematics 84, Springer, New York, 1990.
\bibitem{Se} J.-P. Serre, \emph{A Course in Arithmetic}, Graduate Texts in Mathematics 7, Springer, New York, 1973.
\end{thebibliography}

\begingroup\small
\setlength{\LTcapwidth}{\textwidth}
\begin{longtable}{@{}rllll@{}}
\caption{The cases in Step~3. In each case, the symbol $(A,B)_p$ is computed as the stated symbol, whose entries lie in the stated classes; an entry ``any'' means that the other entry is a unit square, so that Lemma~\ref{lem:square} applies.}\label{tab:local}\\
\hline
 & symbol & first entry & second entry & value\\
\hline\endfirsthead
\hline
 & symbol & first entry & second entry & value\\
\hline\endhead
\hline\endfoot
\multicolumn{5}{@{}l}{$p=2$, solutions modulo $16$}\\[2pt]
1 & $(A,B)$ & $2^{2i+1}\cdot3$ & $2^{2i+1}\cdot3$ & $-1$\\
\hline
\multicolumn{5}{@{}l}{$p=19$, solutions modulo $19$}\\[2pt]
1 & \multicolumn{3}{l}{$19\nmid AB$} & $1$\\
2 & $(A,T_B)$ & $\text{any}$ & $19^{2i}(+)$ & $1$\\
3 & $(A,T_B)$ & $19^{2i}(+)$ & $19^{2i}(+)$ & $1$\\
4 & $(T_A,B)$ & $19^{2i}(-)$ & $19^{2i}(+)$ & $1$\\
\hline
\multicolumn{5}{@{}l}{$p=59$, solutions modulo $59$}\\[2pt]
1 & \multicolumn{3}{l}{$59\nmid AB$} & $1$\\
2 & $(A,T_B)$ & $59^{2i}(+)$ & $59^{2i}(+)$ & $1$\\
3 & $(T_A,B)$ & $59^{2i}(+)$ & $\text{any}$ & $1$\\
4 & $(T_A,B)$ & $59^{2i}(+)$ & $59^{2i}(+)$ & $1$\\
\end{longtable}
\endgroup
\end{document}
